\documentclass[11pt,a4paper]{article}

% Template visual Matriz Aprova. Compile com XeLaTeX ou LuaLaTeX.
\usepackage{fontspec}
\usepackage{polyglossia}
\setmainlanguage{portuguese}
\IfFontExistsTF{Aptos}{
  \setmainfont{Aptos}
  \setsansfont{Aptos}
}{
  \IfFontExistsTF{Calibri}{
    \setmainfont{Calibri}
    \setsansfont{Calibri}
  }{
    \IfFontExistsTF{TeX Gyre Heros}{
      \setmainfont{TeX Gyre Heros}
      \setsansfont{TeX Gyre Heros}
    }{
      \setmainfont{Latin Modern Sans}
      \setsansfont{Latin Modern Sans}
    }
  }
}
\usepackage[a4paper,margin=2cm,headheight=16pt]{geometry}
\usepackage{graphicx}
\usepackage{xcolor}
\usepackage{tikz}
\usepackage{tcolorbox}
\usepackage{enumitem}
\usepackage{booktabs}
\usepackage{tabularx}
\usepackage{amssymb}
\usepackage{fancyhdr}
\usepackage{titlesec}
\usepackage{hyperref}

\definecolor{matrizlime}{HTML}{CBFF4D}
\definecolor{matrizink}{HTML}{101313}
\definecolor{matrizmuted}{HTML}{596363}
\definecolor{matrizline}{HTML}{DDE3DF}
\definecolor{matrizsoft}{HTML}{F3F7EC}

\hypersetup{colorlinks=true,linkcolor=matrizink,urlcolor=matrizink}
\pagestyle{fancy}
\fancyhf{}
\lhead{\textsf{\small MATRIZ APROVA}}
\rhead{\textsf{\small APOSTILA}}
\cfoot{\textsf{\small \thepage}}
\renewcommand{\headrulewidth}{0.4pt}
\renewcommand{\headrule}{\hbox to\headwidth{\color{matrizline}\leaders\hrule height \headrulewidth\hfill}}

\titleformat{\section}{\Large\bfseries\color{matrizink}}{\thesection}{0.6em}{}
\titleformat{\subsection}{\large\bfseries\color{matrizink}}{\thesubsection}{0.6em}{}
\setlist{nosep,leftmargin=*}

\newtcolorbox{foco}{colback=matrizsoft,colframe=matrizlime,boxrule=1pt,arc=3pt,left=10pt,right=10pt,top=8pt,bottom=8pt}
\newtcolorbox{lei}{colback=matrizink,colframe=matrizink,coltext=white,boxrule=0pt,arc=3pt,left=10pt,right=10pt,top=8pt,bottom=8pt}
\newtcolorbox{alerta}{colback=white,colframe=matrizline,boxrule=0.6pt,arc=3pt,left=10pt,right=10pt,top=8pt,bottom=8pt}

% Logo usada na capa. Caminho relativo à pasta onde o .tex é compilado
% (materiais/<disciplina>/). Ajuste se a estrutura mudar.
\newcommand{\logomatriz}{../../public/matrizaprova_temaclaro.png}

\begin{document}

\begin{titlepage}
  \thispagestyle{empty}
  \begin{center}
    \includegraphics[width=0.55\linewidth]{\logomatriz}\par
  \end{center}
  \vspace{2.2cm}
  {\sffamily\fontsize{28}{32}\selectfont\bfseries Apostila — Direito Tributário: Princípios e limitações ao poder de tributar\par}
  \vspace{0.4cm}
  {\sffamily\large Receita Federal, TCU/TCE, CGU e OAB\par}
  \vfill
  \begin{foco}
    \textsf{\textbf{Foco da apostila}}\\[3pt]
    não confunda anterioridade anual com nonagesimal. Um tributo pode respeitar uma e precisar respeitar a outra, conforme a hipótese constitucional.
  \end{foco}
  \vspace{0.7cm}
  {\sffamily\small Atualizado em 16 de agosto de 2026\quad ·\quad Cebraspe, FGV e FCC\par}
  \vspace{1cm}
  {\sffamily\small matrizaprova.com\par}
\end{titlepage}

\tableofcontents
\newpage

% Conteúdo gerado entra aqui.
\section*{O que cai}

O poder de tributar não é ilimitado. A Constituição distribui competências e impõe garantias ao contribuinte. Os pontos mais cobrados são:

\begin{enumerate}
  \item legalidade tributária;
  \item irretroatividade;
  \item anterioridade anual e nonagesimal;
  \item isonomia e capacidade contributiva;
  \item vedação ao confisco;
  \item imunidades e limitações federativas.
\end{enumerate}

\textbf{Regra de prova:} não confunda anterioridade anual com nonagesimal. Um tributo pode respeitar uma e precisar respeitar a outra, conforme a hipótese constitucional.

\section*{Teoria essencial}

\subsection*{1. Poder de tributar e limitações}

Os entes federativos possuem competência tributária definida pela Constituição, mas devem exercê-la dentro das limitações constitucionais. A União, os Estados, o Distrito Federal e os Municípios não podem criar ou aumentar tributos fora de sua competência ou contra as garantias do contribuinte.

As limitações decorrem principalmente de princípios, imunidades e regras de repartição de competências. A lei complementar, o Código Tributário Nacional e as leis ordinárias completam o regime, respeitando a Constituição.

\subsection*{2. Legalidade tributária}

O art. 150, I, veda à União, aos Estados, ao Distrito Federal e aos Municípios exigir ou aumentar tributo sem lei que o estabeleça.

Em regra, a lei deve definir os elementos essenciais da obrigação tributária, como hipótese de incidência, sujeito passivo, base de cálculo e alíquota. A Administração não pode criar tributo por decreto ou ato administrativo.

Há hipóteses constitucionais específicas em que o Poder Executivo pode alterar alíquotas dentro dos limites legais, como ocorre com determinados impostos de função regulatória. Isso não elimina a legalidade: a autorização e os limites devem estar previstos no ordenamento.

\subsection*{3. Isonomia}

É vedado instituir tratamento desigual entre contribuintes que estejam em situação equivalente. A diferenciação pode ser constitucional quando fundada em diferença relevante e relacionada à finalidade da norma.

Isonomia não exige tratar todos de modo matematicamente idêntico. A igualdade tributária permite considerar capacidade econômica, atividade, localização e outras diferenças constitucionalmente relevantes.

\subsection*{4. Irretroatividade}

O art. 150, III, "a", impede cobrar tributos em relação a fatos geradores ocorridos antes do início da vigência da lei que os instituiu ou aumentou.

A regra protege segurança jurídica. A lei tributária nova não pode ser usada para aumentar a carga sobre fato já ocorrido. A análise deve observar a data do fato gerador e a vigência da lei.

Irretroatividade não se confunde com anterioridade: a primeira olha para fatos passados; a segunda impede cobrança imediata de tributo criado ou aumentado.

\subsection*{5. Anterioridade anual}

O art. 150, III, "b", veda cobrar tributo no mesmo exercício financeiro em que haja sido publicada a lei que o instituiu ou aumentou.

Exemplo: se a lei que aumenta um imposto é publicada em maio, a cobrança, em regra, não começa no mesmo exercício; deve aguardar o exercício seguinte, respeitadas outras exigências constitucionais.

Existem exceções previstas na Constituição. Em prova, não memorize a anterioridade como regra absoluta: verifique o tributo e o dispositivo constitucional aplicável.

\subsection*{6. Anterioridade nonagesimal}

O art. 150, III, "c", impede a cobrança antes de decorridos 90 dias da data em que tenha sido publicada a lei que instituiu ou aumentou o tributo.

A contagem é de 90 dias, não necessariamente três meses completos. A regra protege o contribuinte contra aumentos com efeitos imediatos.

Alguns tributos estão excepcionados da noventena ou submetidos a regras especiais. Por isso, questões exigem atenção ao imposto, à contribuição e ao tipo de alteração legislativa.

\subsubsection*{6.1. Combinação das anterioridades}

Quando ambas forem aplicáveis, a cobrança só começa depois de cumpridos os dois marcos: exercício seguinte e 90 dias. Na prática, prevalece a data mais tardia.

\subsection*{7. Capacidade contributiva}

Sempre que possível, os impostos terão caráter pessoal e serão graduados segundo a capacidade econômica do contribuinte, facultado à Administração identificar o patrimônio, os rendimentos e as atividades econômicas, respeitados os direitos individuais e a lei.

Capacidade contributiva busca distribuir o ônus tributário de maneira compatível com a aptidão econômica. Ela não significa que todo tributo precise ser progressivo ou que apenas pessoas com alta renda possam ser tributadas.

\subsection*{8. Não confisco}

É vedado utilizar tributo com efeito de confisco. O tributo não pode ser tão excessivo que, na prática, retire substancialmente o patrimônio ou inviabilize o exercício de atividade lícita.

A análise de confisco considera o caso, a carga global e a finalidade da exação. Não existe um percentual universal que resolva toda questão.

Multas tributárias também podem ser submetidas ao controle de proporcionalidade e à vedação de efeitos confiscatórios, conforme a situação.

\subsection*{9. Liberdade de tráfego}

É vedado estabelecer limitações ao tráfego de pessoas ou bens por meio de tributos interestaduais ou intermunicipais, ressalvada a cobrança de pedágio pela utilização de vias conservadas pelo Poder Público.

O objetivo é impedir que tributos sejam usados como barreiras à circulação interna. O pedágio constitucionalmente admitido possui natureza vinculada ao uso da via, conforme o regime jurídico aplicável.

\subsection*{10. Imunidades tributárias}

Imunidade é uma limitação constitucional ao poder de tributar. A Constituição impede que determinado ente institua imposto em certas situações, pessoas ou objetos.

Não se confunde com isenção:

\begin{tabularx}{\linewidth}{lX}
\toprule
 \textbf{Imunidade} & \textbf{Isenção} \\
\midrule
 decorre da Constituição & decorre de lei \\
 limita a competência tributária & dispensa o pagamento de tributo já previsto \\
 não pode ser livremente revogada por lei ordinária & depende da disciplina legal \\
\bottomrule
\end{tabularx}


Entre as imunidades do art. 150, VI, estão patrimônio, renda ou serviços uns dos outros; templos de qualquer culto; partidos políticos e suas fundações, entidades sindicais dos trabalhadores e instituições de educação e assistência social sem fins lucrativos, observados os requisitos legais; e livros, jornais, periódicos e o papel destinado à impressão.

\subsection*{11. Uniformidade e guerra fiscal}

A União não pode instituir tributo que não seja uniforme em todo o território nacional ou que implique distinção ou preferência em relação a Estado, Distrito Federal ou Município, salvo incentivos destinados a promover o equilíbrio do desenvolvimento socioeconômico entre as diferentes regiões.

Também é vedado aos Estados, ao Distrito Federal e aos Municípios estabelecer diferença tributária entre bens e serviços em razão de sua procedência ou destino.

\section*{Como a banca cobra}

\begin{itemize}
  \item Legalidade não impede todas as alterações de alíquota autorizadas pela Constituição dentro de limites legais.
  \item Irretroatividade trata de fatos geradores passados.
  \item Anterioridade anual trata do exercício financeiro.
  \item Noventena trata do prazo de 90 dias.
  \item Quando ambas são exigidas, vale o prazo mais longo.
  \item Imunidade nasce da Constituição; isenção nasce da lei.
  \item Capacidade contributiva e isonomia não significam tratamento idêntico em situações diferentes.
  \item Pedágio é exceção expressa à liberdade de tráfego.
\end{itemize}

\section*{Exemplos resolvidos}

\subsection*{Exemplo 1}

Uma lei publicada em novembro aumenta determinado imposto. A cobrança pode começar em dezembro do mesmo ano?

\textbf{Em regra, não.} A anterioridade anual impede cobrança no mesmo exercício e, se a noventena também for aplicável, devem transcorrer 90 dias. É necessário verificar as exceções constitucionais do tributo.

\subsection*{Exemplo 2}

Uma lei ordinária concede isenção de imposto a uma categoria. Essa isenção é uma imunidade constitucional?

\textbf{Não.} Imunidade decorre diretamente da Constituição. A dispensa criada por lei é isenção e segue o regime legal correspondente.

\section*{Quadro-resumo}

\begin{tabularx}{\linewidth}{lX}
\toprule
 \textbf{Limitação} & \textbf{Conteúdo} \\
\midrule
 legalidade & tributo exige lei \\
 isonomia & sem tratamento desigual entre equivalentes \\
 irretroatividade & não alcança fato gerador passado \\
 anterioridade anual & não cobra no mesmo exercício da criação/aumento \\
 noventena & aguarda 90 dias \\
 capacidade contributiva & considera aptidão econômica \\
 não confisco & impede carga tributária excessiva \\
 liberdade de tráfego & não cria barreira tributária à circulação \\
 imunidade & limitação constitucional \\
\bottomrule
\end{tabularx}


\section*{Questões de fixação}

\subsection*{Questão 1 — (questão inédita)}

A anterioridade anual impede, em regra, a cobrança de tributo:

A) sobre fato gerador ocorrido antes da lei. \\ B) no mesmo exercício financeiro da publicação da lei que o instituiu ou aumentou. \\ C) depois de decorridos 90 dias. \\ D) sempre que houver alteração de fiscalização. \\ E) apenas de impostos municipais.

\begin{alerta}
\textbf{Gabarito: B.} A hipótese corresponde à anterioridade do exercício financeiro.
\end{alerta}

\subsection*{Questão 2 — (questão inédita)}

A vedação de cobrar tributo antes de decorridos 90 dias da publicação da lei refere-se ao princípio da:

A) isonomia. B) capacidade contributiva. C) noventena. D) legalidade. E) uniformidade.

\begin{alerta}
\textbf{Gabarito: C.} A anterioridade nonagesimal exige o transcurso de 90 dias, ressalvadas as exceções constitucionais.
\end{alerta}

\subsection*{Questão 3 — (questão inédita)}

Assinale a alternativa correta.

A) Imunidade e isenção são sempre sinônimas. \\ B) Imunidade decorre da Constituição, enquanto isenção decorre de lei. \\ C) Isenção sempre decorre de emenda constitucional. \\ D) Imunidade pode ser criada por decreto. \\ E) A isenção define a competência tributária constitucional.

\begin{alerta}
\textbf{Gabarito: B.} A origem constitucional diferencia imunidade de isenção.
\end{alerta}

\subsection*{Questão 4 — (questão inédita)}

O princípio da irretroatividade tributária impede, em regra:

A) cobrança no exercício seguinte. \\ B) cobrança após 90 dias. \\ C) aplicação de lei nova a fato gerador ocorrido antes de sua vigência. \\ D) fiscalização de fatos atuais. \\ E) atualização monetária de valores.

\begin{alerta}
\textbf{Gabarito: C.} Irretroatividade protege fatos geradores anteriores à vigência da lei que instituiu ou aumentou o tributo.
\end{alerta}

\subsection*{Questão 5 — (questão inédita)}

Sobre capacidade contributiva, é correto afirmar que:

A) exige tratamento tributário idêntico para todos. \\ B) impede qualquer tributação de pessoa com baixa renda. \\ C) busca relacionar, quando possível, a carga dos impostos à capacidade econômica. \\ D) aplica-se apenas a taxas. \\ E) autoriza tributo com efeito confiscatório.

\begin{alerta}
\textbf{Gabarito: C.} A capacidade contributiva orienta a graduação dos impostos conforme a aptidão econômica, nos termos constitucionais.
\end{alerta}

\section*{Revisão final}

\begin{itemize}
  \item[$\square$] Sei o conteúdo da legalidade tributária.
  \item[$\square$] Diferencio irretroatividade, anterioridade anual e noventena.
  \item[$\square$] Sei que a data mais tardia prevalece quando anterioridades se combinam.
  \item[$\square$] Diferencio imunidade de isenção.
  \item[$\square$] Entendo capacidade contributiva e isonomia.
  \item[$\square$] Reconheço a vedação ao confisco.
  \item[$\square$] Sei a exceção do pedágio à liberdade de tráfego.
  \item[$\square$] Lembro da vedação de diferença por procedência ou destino.
\end{itemize}

\end{document}
